% =====================================================================
%  CAPÍTULO 6 — CONCLUSIONES
% =====================================================================
\chapter{Conclusiones}
\label{cap:conclusiones}

\section{Conclusiones}
\label{sec:conclusiones}

Este Trabajo Fin de Grado se propuso analizar, mediante técnicas de procesamiento del lenguaje
natural, cómo comunican su estrategia de sostenibilidad las grandes empresas europeas del STOXX
Europe 600 bajo el nuevo marco de la CSRD y los ESRS, y detectar en ese discurso señales
textuales de \textit{greenwashing}. Para ordenar las conclusiones se sigue la recomendación de
la guía de la Facultad y se ofrece una conclusión por cada objetivo específico, de modo que
quede patente que todas las preguntas planteadas al inicio han recibido respuesta.

En relación con el \textbf{primer objetivo} —identificar los temas predominantes—, el análisis
permite concluir que el discurso de sostenibilidad del STOXX 600 se organiza en torno a un
núcleo estable formado por el cambio climático y las personas, con la gobernanza, la economía
circular y la cadena de valor en un segundo plano, y con la contaminación y las comunidades como
los temas menos tratados. Lo relevante no es solo este orden, sino que dos técnicas de modelado
de temas radicalmente distintas —una probabilística y otra basada en el significado— coinciden
en identificar los mismos grandes asuntos. Esa convergencia, reforzada por un tercer clasificador
independiente, otorga una solidez considerable a la radiografía temática del corpus y confirma
que la rejilla de los estándares ESRS captura razonablemente bien la estructura real del
discurso.

En cuanto al \textbf{segundo objetivo} —las diferencias por sector y país—, la conclusión es que
esas diferencias existen y son estadísticamente muy significativas. El sector y la geografía no
son detalles accesorios, sino factores que pesan de forma decisiva en el modo en que se comunica la
sostenibilidad. Destaca, en particular, que los sectores financiero y tecnológico exhiben los
índices de \textit{greenwashing} más altos y, a la vez, la menor especificidad climática,
mientras que los sectores de actividad física intensiva —utilities, inmobiliario, materiales
básicos— reportan con mayor concreción. Este hallazgo conecta directamente con el fenómeno del
\textit{cherry-picking} descrito en la literatura: las empresas con menos que contar en términos
operativos son las que más recurren al lenguaje genérico.

Respecto del \textbf{tercer objetivo} —los determinantes de la especificidad y la relación entre
tono y concreción—, la conclusión obliga a matizar la hipótesis de partida. La idea intuitiva de
que un tono más optimista encubre un discurso menos específico —la señal clásica del
\textit{greenwashing}— no se sostiene en el análisis transversal: ocurre más bien lo contrario,
pues las empresas que hablan de clima con mayor concreción tienden también a un tono más
positivo, probablemente porque la especificidad acompaña a la comunicación de logros reales.
El tamaño de la empresa, por su parte, no resultó un predictor robusto. La señal de
\textit{greenwashing}, por tanto, no debe buscarse en la combinación de optimismo y vaguedad de
una empresa concreta, sino en la pérdida general de sustancia del discurso a lo largo del tiempo,
que es donde el cuarto objetivo aporta la evidencia más clara.

Sobre el \textbf{cuarto objetivo} —la evolución de la NFRD a la CSRD— se obtiene la conclusión
más rotunda del trabajo. El tránsito hacia la obligatoriedad de la CSRD ha transformado
profundamente el reporting de sostenibilidad: los informes casi triplican su extensión, adoptan
un tono marcadamente menos optimista, incorporan mucho más lenguaje de riesgo y de cautela, y
ven crecer su índice de \textit{greenwashing} de forma estadísticamente significativa. El
emblema de ese cambio es la irrupción de la doble materialidad, el concepto nuclear de la CSRD,
cuya presencia en el texto se multiplica por más de nueve en apenas dos años. Ahora bien, el
hallazgo admite una doble lectura que conviene no escamotear. Por un lado, el reporting se vuelve
más extenso pero menos específico y menos cuantitativo, lo que es compatible con un aumento del
\textit{cheap talk}: más palabras que no se traducen en más datos. Por otro, el discurso se
vuelve más prudente, más consciente del riesgo y menos triunfalista, lo que también puede leerse
como una señal de mayor honestidad sobre la incertidumbre, justamente lo que un marco de doble
materialidad pretende fomentar. La evidencia reunida no permite cerrar el dilema, pero sí
afirmar que el primer ejercicio de la CSRD ha producido informes más voluminosos cuyo
crecimiento descansa, en gran medida, sobre contenido descriptivo de cumplimiento más que sobre
nuevos compromisos cuantificados.

De forma transversal, el trabajo permite extraer una conclusión metodológica que trasciende sus
resultados concretos: es posible analizar de manera rigurosa, sistemática y reproducible la
comunicación de sostenibilidad de cientos de empresas sin recurrir a calificaciones ESG de
terceros, interrogando directamente el texto de los informes. La combinación de diccionarios,
modelado de temas y modelos de lenguaje, triangulada con prudencia, ofrece una alternativa
transparente a las cajas negras de los proveedores de \textit{ratings}, y el cuadro de mando
desarrollado demuestra que esa alternativa puede ponerse al servicio de analistas, inversores y
supervisores como herramienta de cribado.

\section{Relación del trabajo con la Agenda 2030}
\label{sec:ods}

La vinculación de este TFG con la Agenda 2030 de Desarrollo Sostenible no es un añadido formal: atañe al núcleo mismo de su objeto. La información de sostenibilidad de las empresas es uno de
los mecanismos a través de los cuales el sector privado rinde cuentas de su contribución a los
Objetivos de Desarrollo Sostenible (ODS), y mejorar la calidad de esa información es, en sí
mismo, una palanca de progreso hacia varios de ellos.

La conexión más directa se establece con el \textbf{ODS 12, de producción y consumo
responsables}, y muy especialmente con su meta 12.6, que insta a alentar a las empresas,
sobre todo a las grandes, a adoptar prácticas sostenibles y a incorporar información sobre
sostenibilidad en su ciclo de presentación de informes. El presente trabajo se sitúa de lleno en
el terreno de esa meta: analiza precisamente la calidad y la evolución de esa información en las
mayores empresas europeas, y aporta herramientas para evaluarla de forma objetiva. Al poner de
manifiesto que el aumento del volumen de los informes no se traduce automáticamente en más
sustancia, contribuye a un debate central para el cumplimiento de la meta 12.6: el de cómo
asegurar que la divulgación sea no solo abundante, sino útil y veraz.

Una segunda conexión, igualmente estrecha, se da con el \textbf{ODS 13, de acción por el clima}.
Buena parte del análisis se concentra en el discurso climático de las empresas —su tono, su
especificidad, su equilibrio entre riesgo y oportunidad—, que es la cara comunicativa de la
meta 13.2 sobre la incorporación de medidas relativas al cambio climático en las estrategias
corporativas. Detectar cuándo el compromiso climático declarado se queda en lo genérico es una
forma de velar por que la acción por el clima sea real y no meramente retórica, lo que entronca
con el espíritu de este objetivo. El trabajo muestra, en este sentido, que el discurso climático
de 2024 se ha vuelto más consciente del riesgo, un cambio de tono que, bien encauzado, puede
favorecer una gestión más seria de la transición.

El trabajo se relaciona también, de manera más indirecta, con el \textbf{ODS 8, de trabajo
decente y crecimiento económico}, en la medida en que la plantilla propia es, junto al clima,
uno de los dos grandes ejes del discurso analizado, y con el \textbf{ODS 16, de paz, justicia e
instituciones sólidas}, cuya meta 16.6 sobre instituciones eficaces y transparentes y cuya
meta 16.10 sobre el acceso público a la información encuentran un reflejo evidente en un trabajo
que promueve la transparencia y la rendición de cuentas corporativas. En última instancia, una
información de sostenibilidad fiable es un bien público que sostiene la confianza en los
mercados y en las instituciones, y este TFG aspira, modestamente, a contribuir a que esa
información sea más escrutable.

\section{Futuras líneas de trabajo}
\label{sec:futuras}

El presente trabajo abre más puertas de las que cierra, y conviene señalar las que parecen más
prometedoras. La primera y más inmediata es la extensión temporal del análisis. Incorporar los
ejercicios posteriores a 2024, a medida que se publiquen, permitiría distinguir si el deterioro
de la especificidad observado en el primer año de la CSRD es un ajuste transitorio —propio del
esfuerzo de adaptación a un marco nuevo— o una tendencia estructural, una pregunta que tres
ejercicios no bastan para responder.

Una segunda línea consistiría en mejorar la precisión de los instrumentos mediante el ajuste
fino de los modelos de lenguaje sobre una muestra de frases de informes ESRS europeos anotada
manualmente. Ello adaptaría las herramientas al dominio específico, permitiría estimar su tasa
de acierto y abriría la puerta a una validación humana sistemática de los temas y de las
clasificaciones, que reforzaría la fiabilidad de los resultados. En la misma dirección apunta la
posibilidad de validar formalmente el índice de \textit{greenwashing} contra casos
independientemente reconocidos, lo que consolidaría su validez de constructo.

Una tercera vía sería ampliar el alcance del corpus más allá de la versión inglesa de los
informes, incorporando el análisis multilingüe de las versiones originales en las distintas
lenguas europeas, lo que evitaría la homogeneización idiomática y permitiría estudiar si el
idioma de redacción influye en el estilo comunicativo. Por último, el cuadro de mando
desarrollado en este trabajo sugiere una línea de carácter más aplicado: convertir el
instrumental analítico en una herramienta operativa de cribado al servicio de supervisores e
inversores, capaz de señalar de forma automática los informes cuyo perfil textual aconseja una
revisión más detenida. El análisis de la comunicación corporativa con inteligencia artificial está, en realidad, dando
sus primeros pasos.
